\chapter{The AIG Managers: GIA and Friends}
\label{ch:aig}

At the core of ABC every design is ultimately an \emph{And-Inverter Graph}
(AIG): a DAG whose only nodes are two-input ANDs, with inversion carried on the
edges. ABC ships several AIG implementations tuned for different jobs. The
modern, scalable one is \pkg{aig/gia} (the \id{Gia_Man_t}, driven by the
\cmd{\&}-commands); the mid-era pointer-based one is \pkg{aig/aig}
(\id{Aig_Man_t}, used by DAR rewriting, DCH choices, and FRAIG); and
\pkg{aig/hop} holds tiny local AIGs used for cut functions inside technology
mapping and rewriting.

\glance{\id{Gia_Man_t} stores nodes in one dense \id{Gia_Obj_t} array; fanins are
\emph{relative} backward offsets (\id{iDiff0}/\id{iDiff1}) and a node is only
two \texttt{unsigned} words plus a \id{Value} word. Literals encode
\texttt{2*var + complement}. Structural hashing (\id{Gia_ManHashAnd}) dedups AND
nodes on creation. Fanouts and levels are optional, computed on demand.}

\section{GIA: the scalable AIG}

\subsection{The node: \id{Gia_Obj_t}}

A GIA object is defined in \file{src/aig/gia/gia.h} as two 32-bit bitfields plus
a scratch word:

\begin{lstlisting}[style=abccode]
struct Gia_Obj_t_
{
    unsigned  iDiff0 : 29;  // the diff of the first fanin
    unsigned  fCompl0:  1;  // the complemented attribute
    unsigned  fMark0 :  1;  // first user-controlled mark
    unsigned  fTerm  :  1;  // terminal node (CI/CO)

    unsigned  iDiff1 : 29;  // the diff of the second fanin
    unsigned  fCompl1:  1;  // the complemented attribute
    unsigned  fMark1 :  1;  // second user-controlled mark
    unsigned  fPhase :  1;  // value under 000 pattern

    unsigned  Value;        // application-specific value
};
\end{lstlisting}

The key trick is that fanins are stored as \emph{relative} offsets, not
absolute IDs or pointers. If a node sits at index \id{Gia_ObjId(p,pObj)}, then
fanin~0 is at index \texttt{id - iDiff0}: literally
\id{Gia_ObjFanin0(pObj) = pObj - pObj->iDiff0}. Because fanins are always
earlier in the array, the diff is a small positive number, the graph is
topologically ordered by construction, and duplicating/renumbering keeps offsets
valid. The whole manager is one dense array (\id{pObjs}), so a node costs only
12 bytes.

The \id{Value} field is scratch space reused by many passes; comments in the
header note it holds the ``next node'' link during structural hashing and the
copy during duplication. Node kind is decoded from the fields, not a type tag:

\begin{table}[htbp]
\centering
\small
\begin{tabular}{@{}lll@{}}
\toprule
Predicate & Condition & Meaning \\
\midrule
\id{Gia_ObjIsConst0} & \id{iDiff0==GIA\_NONE \&\& iDiff1==GIA\_NONE} & object 0, constant \\
\id{Gia_ObjIsCi}     & \id{fTerm \&\& iDiff0==GIA\_NONE}             & CI (PI or LO) \\
\id{Gia_ObjIsCo}     & \id{fTerm \&\& iDiff0!=GIA\_NONE}             & CO (PO or LI) \\
\id{Gia_ObjIsAnd}    & \id{!fTerm \&\& iDiff0!=GIA\_NONE}            & internal AND \\
\id{Gia_ObjIsXor}    & \id{Gia\_ObjIsAnd \&\& iDiff0<iDiff1}         & XOR node \\
\id{Gia_ObjIsBuf}    & \id{iDiff0==iDiff1 \&\& iDiff0!=GIA\_NONE}    & buffer \\
\bottomrule
\end{tabular}
\caption{Node classification is derived from \id{iDiff0}/\id{iDiff1}/\id{fTerm}
(\id{GIA\_NONE} $=$ \texttt{0x1FFFFFFF}, the max 29-bit value).}
\end{table}

\subsection{The manager: \id{Gia_Man_t}}

\id{Gia_Man_t} is a large struct in \file{src/aig/gia/gia.h}; the essential
fields for understanding the data structure are:

\begin{table}[htbp]
\centering
\small
\begin{tabularx}{\linewidth}{@{}llX@{}}
\toprule
Field & Type & Purpose \\
\midrule
\id{pObjs}   & \id{Gia_Obj_t *} & the dense array of all objects \\
\id{nObjs}   & \id{int}         & number of objects in use \\
\id{nObjsAlloc} & \id{int}      & allocated capacity of \id{pObjs} \\
\id{vCis}    & \id{Vec_Int_t *} & IDs of combinational inputs (PIs $+$ LOs) \\
\id{vCos}    & \id{Vec_Int_t *} & IDs of combinational outputs (POs $+$ LIs) \\
\id{nRegs}   & \id{int}         & number of registers (flops) \\
\id{vHTable} & \id{Vec_Int_t}   & structural-hash bucket table \\
\id{vHash}   & \id{Vec_Int_t}   & hash chain links (one entry per object) \\
\id{vRefs} / \id{pRefs} & \id{Vec_Int_t} / \id{int *} & reference (fanout) counts \\
\id{vLevels} & \id{Vec_Int_t *} & per-node logic level \\
\id{pMuxes}  & \id{unsigned *}  & control literals when MUX nodes are present \\
\id{pReprs} / \id{pNexts} & \id{Gia_Rpr_t *} / \id{int *} & equivalence classes \\
\id{pSibls}  & \id{int *}       & choice-node sibling links \\
\id{pFanData} & \id{int *}      & optional explicit fanout database \\
\id{vMapping} / \id{vMapping2} & \id{Vec_Int_t *} / \id{Vec_Wec_t *} & LUT mapping \\
\id{Value} (per obj) & \id{unsigned} & scratch for copies/hash chains \\
\bottomrule
\end{tabularx}
\caption{Selected \id{Gia_Man_t} fields. Most of the remaining $\sim$150 fields
are attachments for specific engines (simulation, timing, abstraction, etc.).}
\end{table}

\subsection{Literals}

GIA never passes bare node pointers between construction calls; it passes
\emph{literals}. A literal packs a variable (object ID) with an inversion bit:
\texttt{lit = 2*var + complement}. The helpers live in
\file{src/misc/util/abc_global.h}:

\begin{table}[htbp]
\centering
\small
\begin{tabular}{@{}ll@{}}
\toprule
Helper & Definition \\
\midrule
\id{Abc_Var2Lit(Var,c)}  & \id{Var + Var + c} \\
\id{Abc_Lit2Var(Lit)}    & \id{Lit >> 1} \\
\id{Abc_LitIsCompl(Lit)} & \id{Lit \& 1} \\
\id{Abc_LitNot(Lit)}     & \id{Lit \^{} 1} \\
\bottomrule
\end{tabular}
\caption{Literal helpers shared across GIA and the SAT layer.}
\end{table}

\subsection{Creating and hashing nodes}

New objects are appended to \id{pObjs} by the inline builders in
\file{src/aig/gia/gia.h}: \id{Gia_ManAppendCi}, \id{Gia_ManAppendAnd},
\id{Gia_ManAppendCo} (and \id{Gia_ManStart}/\id{Gia_ManStop} in
\file{src/aig/gia/giaMan.c} allocate/free the manager). \id{Gia_ManAppendAnd}
computes the two relative diffs, orders fanins, and returns the new node's
literal:

\begin{lstlisting}[style=abccode]
static inline int Gia_ManAppendAnd( Gia_Man_t * p, int iLit0, int iLit1 )
{
    Gia_Obj_t * pObj = Gia_ManAppendObj( p );
    if ( iLit0 < iLit1 ) {
        pObj->iDiff0  = Gia_ObjId(p,pObj) - Abc_Lit2Var(iLit0);
        pObj->fCompl0 = Abc_LitIsCompl(iLit0);
        pObj->iDiff1  = Gia_ObjId(p,pObj) - Abc_Lit2Var(iLit1);
        pObj->fCompl1 = Abc_LitIsCompl(iLit1);
    } else { /* swapped */ }
    return Gia_ObjId( p, pObj ) << 1;
}
\end{lstlisting}

Most passes instead call the \emph{hashed} constructor \id{Gia_ManHashAnd} in
\file{src/aig/gia/giaHash.c}, which folds constants, catches trivial cases, and
looks the pair up in \id{vHTable} before creating a node. On a hit it returns
the existing literal (incrementing \id{nHashHit}); on a miss it calls
\id{Gia_ManAppendAnd}. This structural hashing (``strashing'') guarantees the
AIG stays a maximally-shared canonical two-input form:

\begin{lstlisting}[style=abccode]
int Gia_ManHashAnd( Gia_Man_t * p, int iLit0, int iLit1 )
{
    if ( iLit0 < 2 ) return iLit0 ? iLit1 : 0;   // AND with const
    if ( iLit1 < 2 ) return iLit1 ? iLit0 : 0;
    if ( iLit0 == iLit1 )          return iLit1; // idempotent
    if ( iLit0 == Abc_LitNot(iLit1) ) return 0;  // x & !x = 0
    ...
    int * pPlace = Gia_ManHashFind( p, iLit0, iLit1, -1 );
    if ( *pPlace ) { p->nHashHit++;  return Abc_Var2Lit(*pPlace,0); }
    p->nHashMiss++;
    *pPlace = Abc_Lit2Var( Gia_ManAppendAnd(p, iLit0, iLit1) );
    ...
}
\end{lstlisting}

\subsection{Iterating the graph}

Traversal is done with macros (defined near the end of
\file{src/aig/gia/gia.h}) rather than fanout pointers. Because the array is
topologically ordered, a forward scan visits every node after its fanins:

\begin{table}[htbp]
\centering
\small
\begin{tabularx}{\linewidth}{@{}lX@{}}
\toprule
Macro & Iterates \\
\midrule
\id{Gia_ManForEachObj(p,pObj,i)}  & every object, indices \texttt{0..nObjs-1} \\
\id{Gia_ManForEachAnd(p,pObj,i)}  & internal AND nodes only \\
\id{Gia_ManForEachCi(p,pObj,i)}   & combinational inputs (via \id{vCis}) \\
\id{Gia_ManForEachCo(p,pObj,i)}   & combinational outputs (via \id{vCos}) \\
\id{Gia_ManForEachPi/Po}          & true PIs / POs (excludes registers) \\
\id{Gia_ManForEachRiRo}           & paired register inputs/outputs \\
\bottomrule
\end{tabularx}
\caption{Common GIA iteration macros.}
\end{table}

Fanouts are \emph{not} stored by default. Passes that need them either build the
optional \id{pFanData} database or, more commonly, work with \id{vRefs}
(reference counts) plus \id{vLevels} (logic levels), both computed on demand.
This is the deliberate trade that keeps GIA compact enough to scale to tens of
millions of nodes.

\section{GIA-native synthesis}

The hardcoded \cmd{\&}-scripts live in \file{src/aig/gia/giaScript.c}. The
headline routine is \id{Gia_ManAigSyn2}, which the \cmd{\&syn2} command wraps:
its handler \id{Abc_CommandAbc9Syn2} (in \file{src/base/abci/abc.c}) parses
options and calls
\id{Gia_ManAigSyn2( pAbc->pGia, fOldAlgo, fCoarsen, fCutMin, nRelaxRatio, ... )}.
Internally \cmd{\&syn2} interleaves \id{Gia_ManAreaBalance} with LUT mapping
(\id{Jf_ManPerformMapping}/\id{Lf_ManPerformMapping}) and optional DSD balancing
\-- an AIG-restructuring loop that operates directly on the \id{Gia_Man_t}.

This contrasts with the classic combinational script \cmd{resyn2}, which is a
plain command \emph{alias} (defined in \file{abc.rc}) over the old
network/DAR operators:

\begin{lstlisting}[style=abcshell]
# from abc.rc -- classic AIG rewriting on Abc_Ntk_t / Aig_Man_t
alias resyn2 "b; rw; rf; b; rw; rwz; b; rfz; rwz; b"
# GIA-native equivalent operates in &-space
abc 01> &get; &syn2; &put
\end{lstlisting}

The \cmd{\&get} / \cmd{\&put} pair is the bridge between the two worlds. Their
handlers \id{Abc_CommandAbc9Get} and \id{Abc_CommandAbc9Put} (in
\file{src/base/abci/abc.c}) convert the current \id{Abc_Ntk_t} into
\id{pAbc->pGia} and back, so a design can be pulled into the scalable
\cmd{\&}-engines, optimized, and returned to the classic flow.

\section{Mid-era AIG: \id{Aig_Man_t}}

\file{src/aig/aig/aig.h} defines the pointer-based AIG that predates GIA and is
still used by DAR rewriting (\cmd{dc2}), DCH equivalence/choice computation, and
FRAIG. Nodes carry explicit fanin \emph{pointers} and an object type enum
(\id{AIG_OBJ_AND}, \id{AIG_OBJ_CI}, \dots), inversion is encoded in the low bit
of the pointer:

\begin{lstlisting}[style=abccode]
struct Aig_Obj_t_  // 8 words
{
    union { Aig_Obj_t * pNext; int CioId; };
    Aig_Obj_t *  pFanin0;         // fanin pointers (not offsets)
    Aig_Obj_t *  pFanin1;
    unsigned int Type   : 3;      // explicit object type
    unsigned int fPhase : 1;
    unsigned int fMarkA : 1, fMarkB : 1;
    unsigned int nRefs  : 26;     // reference count
    unsigned     Level  : 24;
    int          TravId, Id;
    union { void * pData; int iData; float dData; };
};
\end{lstlisting}

\begin{table}[htbp]
\centering
\small
\begin{tabularx}{\linewidth}{@{}llX@{}}
\toprule
Field & Type & Purpose \\
\midrule
\id{vCis} / \id{vCos} & \id{Vec_Ptr_t *} & input / output node pointers \\
\id{vObjs}   & \id{Vec_Ptr_t *} & all nodes, indexed by \id{Id} (optional) \\
\id{pConst1} & \id{Aig_Obj_t *} & the constant-1 node \\
\id{pTable}  & \id{Aig_Obj_t **} & structural hash table (chained by \id{pNext}) \\
\id{nObjs[]} & \id{int[]}       & per-type object counts \\
\id{nRegs}   & \id{int}         & number of registers \\
\id{pFanData} / \id{nFansAlloc} & \id{int *} / \id{int} & explicit fanout store \\
\id{pReprs} / \id{pEquivs} & \id{Aig_Obj_t **} & representatives / choice lists \\
\bottomrule
\end{tabularx}
\caption{Selected \id{Aig_Man_t} fields.}
\end{table}

Compared with GIA, \id{Aig_Man_t} trades memory for convenience: objects are
$\sim$8 words, fanins are real pointers, a type tag is stored explicitly, and
fanouts can be materialized. That makes local structural rewriting easy to
express but limits scale, which is exactly why the newer engines were rebuilt on
GIA.

\section{Local AIGs and specialized packages}

\subsection{\id{Hop_Man_t}: tiny function AIGs}

\file{src/aig/hop/hop.h} defines a ``minimalistic'' AIG (\id{Hop_Man_t} /
\id{Hop_Obj_t}). These are small, standalone AIGs used to represent the local
Boolean function of a cut inside technology mapping (\cmd{if}) and rewriting.
The node layout mirrors \id{Aig_Obj_t} (pointer fanins, \id{Type} enum,
complemented-pointer inversion via \id{Hop_Not}/\id{Hop_IsComplement}) but is
lighter (``6 words'') and reuses \id{nRefs} to hold the level.

\subsection{\pkg{aig/saig}: sequential AIGs}

\file{src/aig/saig/saig.h} adds sequential helpers on top of \id{Aig_Man_t},
treating the last \id{nRegs} CIs/COs as latch outputs (LOs) and inputs (LIs).
It provides accessors and iterators such as \id{Saig_ManRegNum},
\id{Saig_ManLo}/\id{Saig_ManLi}, \id{Saig_ObjLoToLi}/\id{Saig_ObjLiToLo}, and
the \id{Saig_ManForEachLiLo} macro, plus register-aware transforms like
\id{Saig_ManDupOrpos} used by the sequential verification engines.

\subsection{Older / niche AIG flavors}

\begin{itemize}
\item \pkg{aig/ivy} \-- an earlier AIG package (\file{src/aig/ivy/ivy.h}) with
its own rewriting; largely superseded by DAR/GIA.
\item \pkg{aig/miniaig} \-- a tiny self-contained reader/writer
(\file{src/aig/miniaig/miniaig.h}, plus \id{minilut}/\id{ndr}) for embedding AIGs
in external tools.
\item \pkg{aig/ioa} \-- AIGER file input/output for the \id{Aig_Man_t} world
(\file{src/aig/ioa/ioa.h}).
\end{itemize}

\section{Choosing a representation}

\begin{table}[htbp]
\centering
\small
\begin{tabularx}{\linewidth}{@{}lXXX@{}}
\toprule
 & \id{Gia_Man_t} & \id{Aig_Man_t} & \id{Abc_Ntk_t} \\
\midrule
Node storage & dense array, relative offsets & heap objects, pointers & heap objects, pointers \\
Node size    & 12 bytes ($2\times$32b $+$ \id{Value}) & $\sim$8 words & larger (names, general gates) \\
Fanins       & \id{iDiff0}/\id{iDiff1} offsets & \id{pFanin0}/\id{pFanin1} & fanin vector \\
Fanouts      & optional (\id{pFanData}) / refs & optional / explicit & always (fanin+fanout vectors) \\
Node kinds   & 2-input AND (XOR/MUX/buf) & AND/EXOR $+$ type enum & arbitrary gates, LUTs, SOPs \\
Scale        & tens of millions & millions & smaller, logic-level \\
Typical use  & \cmd{\&}-flow: mapping, \cmd{\&dch}, \cmd{\&syn2} & DAR \cmd{dc2}, DCH, FRAIG & user-facing netlist, I/O, general logic \\
\bottomrule
\end{tabularx}
\caption{GIA vs.\ \id{Aig_Man_t} vs.\ \id{Abc_Ntk_t}.}
\end{table}

The practical rule of thumb: \id{Abc_Ntk_t} is the general-purpose netlist the
user sees; \id{Aig_Man_t} is the working AIG for pointer-friendly rewriting and
choice-based engines; and \id{Gia_Man_t} is the compact, scalable AIG behind the
\cmd{\&}-commands where memory footprint and raw size matter most. \cmd{\&get}
and \cmd{\&put} move designs between the \id{Abc_Ntk_t} and \id{Gia_Man_t}
worlds.
